\section{The Method: Two Integers and One Trunk Choice}\label{sec:method}

This section states the entire method before any variant is named. Two plain-language terms are
used throughout. The \emph{trunk} of a block is the stack of fully connected layers that maps
the input window to a hidden vector. A \emph{basis} is a fixed set of orthonormal vectors into
which a block's output is expanded.

Each block of a doubly residual stack \citep{oreshkin2020nbeats} carries two unconstrained
linear synthesis heads, one producing a backcast of length $\L$ and one a forecast of length
$\Ho$. The method replaces each head by a factorisation through a fixed orthonormal subspace.
With $\mathbf{x}$ the block input and $\mathbf{h}$ the trunk output,
\begin{align}
\mathbf{h} &= \operatorname{trunk}(\mathbf{x}),
  &&\text{trunk: uniform MLP or bottleneck,}
  \label{eq:trunk}\\
\boldsymbol{\theta}^{b} &= \mathbf{W}_{b}\,\mathbf{h},
  &\boldsymbol{\theta}^{f} &= \mathbf{W}_{f}\,\mathbf{h},
  &&\text{learned, unconstrained,}
  \label{eq:heads}\\
\hat{\mathbf{x}} &= \Bb^{\top}\boldsymbol{\theta}^{b},
  &\hat{\mathbf{y}} &= \Bf^{\top}\boldsymbol{\theta}^{f},
  &&\Bb,\ \Bf\ \text{frozen, orthonormal rows.}
  \label{eq:synthesis}
\end{align}
The projections $\mathbf{W}_{b}$ and $\mathbf{W}_{f}$ are learned and unconstrained; the basis
matrices $\Bb \in \mathbb{R}^{\db \times \L}$ and $\Bf \in \mathbb{R}^{\df \times \Ho}$ are
frozen at build time with orthonormal rows. The stack remains doubly residual: the backcast
$\hat{\mathbf{x}}$ is subtracted from the block input and the forecast $\hat{\mathbf{y}}$ is
added to the running output. The method therefore reduces to two integers, the ranks $\db$ and
$\df$, and one trunk choice, uniform width or bottleneck. Every named variant in this paper is a
choice of how the rows of $\Bb$ and $\Bf$ are generated or how $\mathbf{h}$ is produced;
optional branches add a frozen polynomial trend basis of order $p$ or a learned generic branch
of rank $r$ to the same additive decomposition.

Table~\ref{tab:dims} fixes the notation. Every symbol used in the paper is introduced exactly
once, in that table; nothing outside it introduces a symbol.

\begin{table}[t]
\centering
\small
\caption{Unified notation and dimensions. Every symbol in the paper appears exactly once here.
Defaults used unless stated: $\S = 10$ stacks, $\U = 512$, $p = 3$; the forecast basis is
complete when $\df = \Ho$.}
\label{tab:dims}
\begin{tabular}{@{}lll@{}}
\toprule
Symbol & Meaning & Dimension \\
\midrule
$\L$ & backcast (input window) length & $\mathbb{N}$ \\
$\Ho$ & forecast (horizon) length & $\mathbb{N}$ \\
$\S$ & number of stacks & $\mathbb{N}$ \\
$\U$ & uniform-trunk hidden width & $\mathbb{N}$ \\
$\U/2$ & bottleneck-trunk half width & $\mathbb{N}$ \\
$d$ & bottleneck latent width & $\mathbb{N}$ \\
$\db$ & backcast basis rank & $\mathbb{N}$, $\db \le \L$ \\
$\df$ & forecast basis rank & $\mathbb{N}$, $\df \le \Ho$ \\
$p$ & polynomial order of the trend branch & $\mathbb{N}$ \\
$r$ & learned generic-branch rank & $\mathbb{N}$ \\
$\mathbf{x}$ & block input (residual stream) & $\mathbb{R}^{\L}$ \\
$\mathbf{h}$ & trunk output & $\mathbb{R}^{\U}$ \\
$\mathbf{W}_{b}$ & learned backcast projection & $\mathbb{R}^{\db \times \U}$ \\
$\mathbf{W}_{f}$ & learned forecast projection & $\mathbb{R}^{\df \times \U}$ \\
$\boldsymbol{\theta}^{b}$ & backcast coefficient vector & $\mathbb{R}^{\db}$ \\
$\boldsymbol{\theta}^{f}$ & forecast coefficient vector & $\mathbb{R}^{\df}$ \\
$\Bb$ & frozen backcast basis, orthonormal rows & $\mathbb{R}^{\db \times \L}$ \\
$\Bf$ & frozen forecast basis, orthonormal rows & $\mathbb{R}^{\df \times \Ho}$ \\
$\hat{\mathbf{x}}$ & backcast output & $\mathbb{R}^{\L}$ \\
$\hat{\mathbf{y}}$ & forecast output & $\mathbb{R}^{\Ho}$ \\
\bottomrule
\end{tabular}
\end{table}

Table~\ref{tab:taxonomy} fixes the variant taxonomy. Variant names carry branch composition
only; depth, trunk, ranks, layout, and offsets are table columns, never name tokens. The
wavelet-family field leaves every configuration string; the axis leaves the notation and roughly
a third of the tables.

\begin{table}[t]
\centering
\small
\caption{Variant taxonomy. Variant names encode branch composition only; every other axis
(depth $\S$, ranks $\db$/$\df$, trunk width, layout, offsets) is a table column. N-BEATS-G and
N-BEATS-I+G are the baselines of \citet{oreshkin2020nbeats} under their original names.}
\label{tab:taxonomy}
\begin{tabular}{@{}llll@{}}
\toprule
Variant & Trunk & Forecast head & Branches \\
\midrule
N-BEATS-G   & uniform             & unconstrained linear                        & none \\
N-BEATS-I+G & uniform             & basis expansion (polynomial + Fourier)      & trend, seasonality \\
FB          & uniform or bottleneck & frozen orthonormal basis                    & none \\
FB+T        & uniform or bottleneck & frozen orthonormal basis                    & polynomial trend (order $p$) \\
FB+T+G      & uniform or bottleneck & frozen orthonormal basis                    & trend + learned generic (rank $r$) \\
NHiTS-FB    & NHiTS backbone \citep{challu2023nhits} & frozen orthonormal basis & none \\
\bottomrule
\end{tabular}
\end{table}

Three markers recur in the results tables; they are defined once, in
Table~\ref{tab:legend}, and nowhere else.

\begin{table}[t]
\centering
\small
\caption{Unified legend for the markers used in all results tables. Parameter counts throughout
the paper include the frozen bases.}
\label{tab:legend}
\begin{tabular}{@{}lp{0.72\columnwidth}@{}}
\toprule
Marker & Meaning \\
\midrule
$\dagger$ & \texttt{active\_g} on the forecast head: the trunk activation is applied to the
forecast head. A \texttt{lightningnbeats} extension, not part of the original N-BEATS paper
\citep{oreshkin2020nbeats}; ``paper-faithful'' claims exclude $\dagger$ rows. \\
\textsuperscript{g} & learned-gate bottleneck variant; merged into the bottleneck pool in
accuracy tables because the gate is null (minimum $p = 0.24$ over 12 comparisons). \\
$\ddagger$ & identity-basis cell: the wavelet generator returned the identity matrix because the
horizon is shorter than the filter's minimum length; the label ``wavelet'' does not apply to
this cell. \\
\bottomrule
\end{tabular}
\end{table}

The bottleneck trunk replaces the uniform hidden width by the width sequence $\U \to \U/2 \to d
\to \U/2 \to \U$; an optional learned gate rescales the latent vector after the latent layer,
\begin{equation}
\mathbf{z} \leftarrow \mathrm{sigm}(\mathbf{v}) \odot \mathbf{z},
\label{eq:gate}
\end{equation}
with $\mathbf{v} \in \mathbb{R}^{d}$ learned and $\odot$ elementwise multiplication. This is the
only place the gate appears. It is null in practice --- twelve paired comparisons, minimum
$p = 0.24$ --- so accuracy tables pool gated and ungated bottleneck rows and flag the gated ones
with the \textsuperscript{g} marker (Table~\ref{tab:legend}).

The scope of the results built on this formulation is fixed by four conditions, stated here once
and recalled wherever they bind. First, the completeness argument of
Section~\ref{sec:completeness} covers the forecast path only: every shipped backcast basis is
truncated to $\Ho$ or $2\Ho$ rows against a backcast of length $\L$ between $3\Ho$ and $5\Ho$,
and truncated bases from different families span different subspaces that no rotation relates.
Second, tiered cells truncate the forecast basis and belong to the truncated regime of
Section~\ref{sec:truncation}, not to the complete-rank regime. Third, every arm in the
identity-swap grid carried a paired trend branch, so the measured null is a
trend-branch-present null. Fourth, the guarantee is equal expressiveness, not equal
trainability; Section~\ref{sec:trainability} shows the two come apart.

The remainder of the paper is organized around the four results this formulation supports.
Section~\ref{sec:trunk} shows that the parameters are in the trunk and are not needed.
Section~\ref{sec:completeness} proves that basis identity is inert at complete forecast rank and
measures the boundary of that guarantee. Section~\ref{sec:truncation} shows that captured
energy predicts accuracy once the basis is truncated. Section~\ref{sec:trainability} shows that
effective rank predicts whether the block trains at all, and gives the constructive repairs.

